\documentclass[10pt]{article}
\usepackage[margin=0.75in]{geometry}
\usepackage{amsmath, amssymb, amsthm}
\usepackage{microtype}
\usepackage{enumitem}
\usepackage{fancyhdr}
\usepackage{titlesec}
\usepackage{tcolorbox}
\usepackage{booktabs}

\linespread{1.08}
\setlength{\parskip}{0.4ex plus 0.2ex minus 0.1ex}

\titleformat{\section}{\large\bfseries}{\thesection.}{0.5em}{}
\titleformat{\subsection}{\normalsize\bfseries}{\thesubsection}{0.5em}{}
\titlespacing*{\section}{0pt}{1.5ex}{1ex}
\titlespacing*{\subsection}{0pt}{1ex}{0.5ex}

\setlist{itemsep=1pt, topsep=3pt, parsep=1pt, leftmargin=1.5em}

\pagestyle{fancy}
\fancyhf{}
\fancyhead[L]{\small EECS 127}
\fancyhead[R]{\small Discussion 3}
\fancyfoot[C]{\small\thepage}
\renewcommand{\headrulewidth}{0.4pt}

\newcommand{\RR}{\mathbb{R}}
\newcommand{\NN}{\mathcal{N}}
\newcommand{\Range}{\mathcal{R}}
\DeclareMathOperator{\rank}{rank}
\DeclareMathOperator{\trace}{tr}
\DeclareMathOperator*{\argmax}{arg\,max}
\DeclareMathOperator*{\argmin}{arg\,min}

\begin{document}

\begin{center}
{\large\bfseries EECS 127/227AT \hfill Optimization Models in Engineering \hfill UC Berkeley \hfill Spring 2026}\\[6pt]
{\LARGE\bfseries Discussion 3}
\end{center}

\section{SVD}

Suppose we have a matrix $A \in \RR^{m \times n}$ with rank $r$. It turns out that its SVD has multiple forms, all of which can be useful depending on the problem we're working on.

We define the \underline{compact SVD} as follows:
\[
\underbrace{A}_{m \times n} = \underbrace{U_r}_{m \times r} \underbrace{\Sigma_r}_{r \times r} \underbrace{V_r^\top}_{r \times n}.
\]

Here, $\Sigma_r \in \RR^{r \times r}$ is a diagonal matrix containing non-zero singular values of $A$.
\[
\Sigma_r = \begin{bmatrix} \sigma_1 & & \\ & \ddots & \\ & & \sigma_r \end{bmatrix},
\]
with $\sigma_1 \geq \sigma_2 \geq \ldots \geq \sigma_r$.

Next, $U_r \in \RR^{m \times r}$ is given by,
\[
U_r = \begin{bmatrix} \vec{u}_1, \vec{u}_2, \ldots, \vec{u}_r \end{bmatrix},
\]
where $u_i$ is a left singular vector corresponding to non-zero singular value, $\sigma_i$, for $i = 1, 2, \ldots, r$. The columns of $U_r$ are orthonormal and together they span the columnspace of $A$.

Finally, $V_r^\top \in \RR^{r \times n}$ is given by,
\[
V_r^\top = \begin{bmatrix} \vec{v}_1^\top \\ \vec{v}_2^\top \\ \vdots \\ \vec{v}_r^\top \end{bmatrix},
\]
where $\vec{v}_j$ is a right singular vector corresponding to non-zero singular value, $\sigma_j$ for $j = 1, 2, \ldots, r$. The rows of $V_r^\top$ are orthonormal and span the rowspace of $A$. Equivalently the columns of $V_r$ span the column space of $A^\top$.

The matrix $A$ can be expressed as,
\[
A = \sigma_1 \vec{u}_1 \vec{v}_1^\top + \sigma_2 \vec{u}_2 \vec{v}_2^\top + \ldots + \sigma_r \vec{u}_r \vec{v}_r^\top.
\]
This is called the \underline{dyadic SVD}, since it's expressed as the sum of dyads (matrices of the form $uv^\top$). Assume now that $m \geq n$.

Another type of SVD which might be more familiar is the \underline{full SVD} of $A$ which is defined as follows:
\[
\underbrace{A}_{m \times n} = \underbrace{U}_{m \times m} \underbrace{\Sigma}_{m \times n} \underbrace{V^\top}_{n \times n}.
\]

Here, $\Sigma \in \RR^{m \times n}$ has non-diagonal entries as zero. The diagonal entries of $\Sigma$ contain the singular values and we can write $\Sigma$ in terms of $\Sigma_r$ as,
\[
\Sigma = \begin{bmatrix} \Sigma_r & 0_{r \times (n-r)} \\ 0_{(m-r) \times r} & 0_{(m-r) \times (n-r)} \end{bmatrix}
\]

Next, $U \in \RR^{m \times m}$ is an orthonormal matrix. $U$ can be expressed in terms of $U_r$ as,
\[
U = \begin{bmatrix} \underbrace{U_r}_{m \times r} & \underbrace{\vec{u}_{r+1} \quad \ldots \quad \vec{u}_m}_{m \times (m-r)} \end{bmatrix}
\]
The columns $\vec{u}_{r+1}, \vec{u}_{r+2}, \ldots, \vec{u}_m$ are left singular vectors corresponding to singular value $0$, and together span the nullspace of $A^\top$.

Finally, $V^\top$ is an orthonormal matrix and can be expressed in terms of $V_r^\top$ as,
\[
V^\top = \begin{bmatrix} V_r^\top \\ \vec{v}_{r+1}^\top \\ \vdots \\ \vec{v}_n^\top \end{bmatrix}
\left.
\begin{array}{l}
\left.\vphantom{\begin{matrix} V_r^\top \end{matrix}}\right\} r \times n \\[6pt]
\left.\vphantom{\begin{matrix} \vec{v}_{r+1}^\top \\ \vdots \\ \vec{v}_n^\top \end{matrix}}\right\} (n-r) \times n
\end{array}
\right.
\]
The rows $\vec{v}_{r+1}^\top, \vec{v}_{r+2}^\top, \ldots, \vec{v}_n^\top$ when transposed are the right singular vectors corresponding to singular value $0$, and together they span the nullspace of $A$.

\begin{enumerate}[label=(\alph*)]
\item For this problem assume that $m > n > r$. Label each of the following as True or False:

\begin{enumerate}[label=(\alph*)]
    \item $UU^\top = I$

    \vspace{3cm}

    \item $U^\top U = I$

    \vspace{3cm}

    \item $V^\top V = I$

    \vspace{3cm}

    \item $VV^\top = I$

    \vspace{3cm}

    \item $U_r^\top U_r = I$

    \vspace{3cm}

    \item $U_r U_r^\top = I$

    \vspace{3cm}

    \item $V_r V_r^\top = I$

    \vspace{3cm}

    \item $V_r^\top V_r = I$

    \vspace{3cm}
\end{enumerate}

\item Find the compact SVD of $A$, given that it has the following full SVD:
\[
A = \begin{bmatrix} 0 & 1 & 0 \\ 1 & 0 & 0 \\ 0 & 0 & 1 \end{bmatrix} \begin{bmatrix} 2 & 0 \\ 0 & 0 \\ 0 & 0 \end{bmatrix} \begin{bmatrix} 1 & 0 \\ 0 & 1 \end{bmatrix}.
\]

\vspace{6cm}

\item Find the full SVD of $A$, given that it has the following compact SVD:
\[
A = \begin{bmatrix} \frac{1}{\sqrt{2}} & 0 \\[4pt] \frac{1}{\sqrt{2}} & 0 \\[4pt] 0 & 1 \end{bmatrix} \begin{bmatrix} 2 & 0 \\ 0 & 1 \end{bmatrix} \begin{bmatrix} 1 & 0 \\ 0 & 1 \end{bmatrix}.
\]

\vspace{6cm}

\end{enumerate}

\section{PCA and Regression}

\begin{enumerate}[label=(\alph*)]
\item Given the following plot of data in $\RR^2$ (i.e., each dot is a data point in $\RR^2$) and candidate unit vectors $\vec{v}_1, \vec{v}_2, \vec{v}_3, \vec{v}_4 \in \RR^2$, \textbf{identify the candidate vectors which could be the first principal component and second principal component of the data (and specify which is which).}

\begin{center}
\emph{[Scatter plot with data points forming an elongated cloud tilted upper-right, with four candidate unit vectors $\vec{v}_1$ (pointing left), $\vec{v}_2$ (pointing upper-left), $\vec{v}_3$ (pointing upper-right along the data spread), and $\vec{v}_4$ (pointing down).]}
\end{center}

\vspace{6cm}

\item Suppose we have data $\vec{x}_1, \ldots, \vec{x}_n \in \RR^d$, where $n > d$. We arrange these data points into a matrix, i.e.,
\begin{equation}
X = \begin{bmatrix} \vec{x}_1^\top \\ \vdots \\ \vec{x}_n^\top \end{bmatrix} \in \RR^{n \times d}.
\end{equation}
\emph{Assume that $X$ is centered, i.e., each column has mean zero:} $(1/n) \sum_{i=1}^{n} \vec{x}_i = \vec{0}_d$, \emph{where $\vec{0}_d$ is the zero vector in $\RR^d$.} Suppose that $X$ has compact SVD given by $X = U_d \Sigma_d V_d^\top$ where
\begin{equation}
U_d = \begin{bmatrix} \vec{u}_1, \ldots, \vec{u}_d \end{bmatrix} \in \RR^{n \times d}, \qquad V_d = \begin{bmatrix} \vec{v}_1, \ldots, \vec{v}_d \end{bmatrix} \in \RR^{d \times d}, \qquad \Sigma_d = \begin{bmatrix} \sigma_1 & & \\ & \ddots & \\ & & \sigma_d \end{bmatrix} \in \RR^{d \times d}
\end{equation}
where $\sigma_1 > \sigma_2 > \cdots > \sigma_d > 0$. \textbf{From this SVD, identify the top $k$ principal components of the data $\{\vec{x}_1, \ldots, \vec{x}_n\} \subseteq \RR^d$, where $k \leq d$.}

\emph{HINT: Recall that the first principal component solves the optimization problem} $\displaystyle\argmax_{\vec{w} \in \RR^d :\, \|\vec{w}\|_2 = 1} \vec{w}^\top X^\top X \vec{w}$.

\vspace{6cm}

\end{enumerate}

\section{Singular Values}

\begin{enumerate}[label=(\alph*)]
\item Suppose $A \in \RR^{3 \times 2}$ is a matrix such that $A^\top A$ is given by
\begin{equation}
A^\top A = \begin{bmatrix} 1/\sqrt{2} & -1/\sqrt{2} \\ 1/\sqrt{2} & 1/\sqrt{2} \end{bmatrix} \begin{bmatrix} 5 & 0 \\ 0 & 3 \end{bmatrix} \begin{bmatrix} 1/\sqrt{2} & 1/\sqrt{2} \\ -1/\sqrt{2} & 1/\sqrt{2} \end{bmatrix}.
\end{equation}
\textbf{What are the singular values of $A$?}

\vspace{6cm}

\item Suppose that $B \in \RR^{3 \times 2}$ has singular values $\sqrt{2}$ and $\sqrt{7}$. Let $C = \begin{bmatrix} B & -B & 3I_3 \end{bmatrix} \in \RR^{3 \times 7}$, where $I_3 \in \RR^{3 \times 3}$ is the $3 \times 3$ identity matrix. \textbf{What are the singular values of $C$?}

\emph{HINT: Consider the matrix $CC^\top \in \RR^{3 \times 3}$.}

\vspace{6cm}

\end{enumerate}

\end{document}
